\documentclass[assignment.tex]{subfiles}
\begin{document}

\chapter{Assignment 2: Semantic Extraction}
\label{sec:assignment2}

Building on the syntactic pipeline of Assignment~1, this stage extracts
\emph{semantic} information from each clause: named entities that identify the
parties, values, and legal references involved; semantic roles that describe who
does what to whom; and an intent label that captures the functional type of the
clause.

\section{Custom Named Entity Recognition (NER)}
\label{subsec:ner}

\subsection{Entity Schema}

Six domain-specific entity types are defined for legal contract NER:

\begin{itemize}[leftmargin=*, itemsep=3pt]
    \item \textbf{PARTY} — contracting parties and organisations
          (e.g.\ \textit{Party A}, \textit{ABC International Co., Ltd.}).
    \item \textbf{MONEY} — monetary amounts with or without currency symbols
          (e.g.\ \textit{USD 10,000}, \textit{\$5,000}).
    \item \textbf{DATE} — absolute and relative date expressions
          (e.g.\ \textit{05 May 2024}, \textit{within 30 days}).
    \item \textbf{RATE} — percentage or per-unit rates
          (e.g.\ \textit{1\% per month}, \textit{2\% per annum}).
    \item \textbf{PENALTY} — penalty clauses and liquidated-damages expressions.
    \item \textbf{LAW} — statutory references, ordinances, and legal codes
          (e.g.\ \textit{Civil Code 2015}).
\end{itemize}

The full BIO label set used during model training is:
\begin{center}
    \texttt{O, B-PARTY, I-PARTY, B-MONEY, I-MONEY, B-DATE, I-DATE,
    B-RATE, I-RATE, B-PENALTY, I-PENALTY, B-LAW, I-LAW}
\end{center}

\subsection{Model Architecture}

\paragraph{Primary model — Legal-BERT fine-tuned NER.}
The primary model fine-tunes
\texttt{nlpaueb/legal-bert-base-uncased}~\cite{chalkidis2020legal}
for token classification using the HuggingFace \texttt{Trainer} API.
Key training hyperparameters are summarised in
Table~\ref{tab:ner-hparams}.

\begin{table}[H]
    \centering
    \caption{Legal-BERT NER training configuration.}
    \label{tab:ner-hparams}
    \begin{tabular}{@{}ll@{}}
        \toprule
        \textbf{Hyperparameter} & \textbf{Value} \\
        \midrule
        Base model           & \texttt{nlpaueb/legal-bert-base-uncased} \\
        Epochs               & 3 \\
        Learning rate        & $2 \times 10^{-5}$ \\
        Batch size           & 8 \\
        Train/eval split     & 90\,/\,10 \\
        Evaluation metric    & seqeval (entity-level F1) \\
        \bottomrule
    \end{tabular}
\end{table}

\paragraph{Fallback — Rule-based NER.}
When Legal-BERT confidence falls below a configurable threshold, a
\texttt{RuleBasedNER} module is invoked.  It applies hand-crafted regular
expressions tailored to each entity type — for example, monetary patterns of the
form \texttt{(USD|VND|\textbackslash\$)\textbackslash s*[\textbackslash
d,]+} for MONEY, and ISO-format or natural-language date patterns for DATE.

\subsection{Output Format}

Results are written to \texttt{output/ner\_results.json} as a list of objects:

\begin{lstlisting}[style=jsonStyle,
                   caption={NER output format (abbreviated).},
                   label={lst:ner-json}]
[
  {
    "clause_id": 0,
    "text": "Party A shall pay Party B USD 10,000 by 05 May 2024...",
    "entities": [
      {"text": "Party A",     "label": "PARTY",  "start":  0, "end":  7},
      {"text": "Party B",     "label": "PARTY",  "start": 16, "end": 23},
      {"text": "USD 10,000",  "label": "MONEY",  "start": 24, "end": 34},
      {"text": "05 May 2024", "label": "DATE",   "start": 38, "end": 49},
      {"text": "1% per month","label": "RATE",   "start": 53, "end": 65}
    ]
  }
]
\end{lstlisting}

\subsection{NER Example}

Table~\ref{tab:ner-example} shows the entities extracted from a representative
legal clause.

\begin{table}[H]
    \centering
    \caption{Entities extracted from \textit{``Party A shall pay Party B USD 10,000
    by 05 May 2024 at 1\% per month penalty rate.''}}
    \label{tab:ner-example}
    \begin{tabular}{@{}llrr@{}}
        \toprule
        \textbf{Entity Text} & \textbf{Label} & \textbf{Start} & \textbf{End} \\
        \midrule
        Party A      & PARTY  &  0 &  7 \\
        Party B      & PARTY  & 16 & 23 \\
        USD 10,000   & MONEY  & 24 & 34 \\
        05 May 2024  & DATE   & 38 & 49 \\
        1\% per month & RATE  & 53 & 65 \\
        \bottomrule
    \end{tabular}
\end{table}

\section{Semantic Role Labelling (SRL)}
\label{subsec:srl}

\subsection{Approach}

Rather than a dedicated BERT-based SRL model, the system uses a
\emph{dependency-tree-based} approach built on top of spaCy's
\texttt{en\_core\_web\_sm} parser.  For each clause, the \texttt{ROOT} token
(typically the main verb) is identified as the predicate, and semantic roles are
assigned to its dependents by mapping Universal Dependencies labels to
PropBank-style argument roles, as shown in Table~\ref{tab:srl-mapping}.

\begin{table}[H]
    \centering
    \caption{Mapping from Universal Dependencies labels to semantic roles.}
    \label{tab:srl-mapping}
    \begin{tabular}{@{}lll@{}}
        \toprule
        \textbf{UD Dependency Label} & \textbf{Semantic Role} & \textbf{Notes} \\
        \midrule
        \texttt{nsubj}              & Agent      & Active-voice subject \\
        \texttt{nsubjpass}          & Theme      & Passive-voice subject \\
        \texttt{dobj} / \texttt{obj}& Theme      & Direct object \\
        \texttt{iobj}               & Recipient  & Indirect object \\
        \texttt{prep} (to-)         & Recipient  & Prepositional recipient \\
        Temporal prepositions       & Time       & \texttt{prep} with time head \\
        \texttt{advcl} + cond.\ mark& Condition  & Subordinate conditional clause \\
        \bottomrule
    \end{tabular}
\end{table}

\subsection{Output Format}

Results are written to \texttt{output/srl\_results.json} as a list of objects with
the following structure:

\begin{lstlisting}[style=jsonStyle,
                   caption={SRL output format.},
                   label={lst:srl-json}]
[
  {
    "clause_id": 0,
    "text": "The Employer shall pay the Employee the monthly salary ...",
    "predicate": "pay",
    "roles": {
      "Agent":     "The Employer",
      "Theme":     "the monthly salary",
      "Recipient": "the Employee",
      "Time":      "within 30 days"
    }
  }
]
\end{lstlisting}

\subsection{SRL Example}

\noindent\textbf{Input:}
\begin{quote}
    \textit{``The Employer shall pay the Employee the monthly salary within 30 days.''}
\end{quote}

\noindent\textbf{Output:}
\begin{quote}
\begin{tabular}{@{}ll@{}}
    Predicate: & \textit{pay} \\
    Agent:     & \textit{The Employer} \\
    Theme:     & \textit{the monthly salary} \\
    Recipient: & \textit{the Employee} \\
    Time:      & \textit{within 30 days} \\
\end{tabular}
\end{quote}

The ROOT token \textit{pay} is the predicate.  \textit{The Employer} is the active
nominal subject (\texttt{nsubj}) mapped to Agent; \textit{the monthly salary} is the
direct object (\texttt{dobj}) mapped to Theme; \textit{the Employee} is reached via
the indirect object (\texttt{iobj}) relation mapped to Recipient; and
\textit{within 30 days} is a temporal prepositional phrase mapped to Time.

\section{Intent Classification}
\label{subsec:intent}

\subsection{Task Definition and Label Set}

Intent classification assigns each clause one of four functional intent labels:

\begin{itemize}[leftmargin=*, itemsep=2pt]
    \item \textbf{Obligation} — a party is required to perform an action
          (keywords: \textit{shall, must, is obligated to}).
    \item \textbf{Prohibition} — an action is forbidden
          (keywords: \textit{shall not, must not, is prohibited}).
    \item \textbf{Right} — a party is entitled to perform an action
          (keywords: \textit{may, has the right to, is entitled to}).
    \item \textbf{Termination Condition} — the contract terminates upon an event
          (keywords: \textit{terminates, expires, is terminated upon/if}).
\end{itemize}

A priority order — Prohibition $>$ Termination Condition $>$ Right $>$ Obligation
— is used to resolve clauses matching multiple keyword patterns during weak-label
generation.

\subsection{Model 1 — TF-IDF + Logistic Regression}

Training labels are generated by a keyword-regex weak labeller.  A TF-IDF
vectoriser converts each clause to a feature vector, and a Logistic Regression
classifier is trained on these vectors.  Table~\ref{tab:tfidf-config} summarises
the configuration.

\begin{table}[H]
    \centering
    \caption{TF-IDF + Logistic Regression configuration.}
    \label{tab:tfidf-config}
    \begin{tabular}{@{}ll@{}}
        \toprule
        \textbf{Component} & \textbf{Setting} \\
        \midrule
        \multicolumn{2}{@{}l}{\textit{TfidfVectorizer}} \\
        \quad n-gram range   & $(1, 3)$ \\
        \quad max features   & 5\,000 \\
        \quad sublinear TF   & \texttt{True} \\[4pt]
        \multicolumn{2}{@{}l}{\textit{LogisticRegression}} \\
        \quad Regularisation $C$  & 1.0 \\
        \quad Solver         & \texttt{lbfgs} \\
        \quad Multi-class    & \texttt{multinomial} \\
        \quad Max iterations & 1\,000 \\
        \bottomrule
    \end{tabular}
\end{table}

\subsection{Model 2 — Fine-tuned BERT}

The second classifier fine-tunes \texttt{bert-base-uncased} as a sequence
classifier using \texttt{AutoModelForSequenceClassification}.
Table~\ref{tab:bert-intent-config} summarises training settings.

\begin{table}[H]
    \centering
    \caption{BERT intent classification training configuration.}
    \label{tab:bert-intent-config}
    \begin{tabular}{@{}ll@{}}
        \toprule
        \textbf{Hyperparameter} & \textbf{Value} \\
        \midrule
        Base model       & \texttt{bert-base-uncased} \\
        Epochs           & 3 \\
        Learning rate    & $2 \times 10^{-5}$ \\
        Max token length & 256 \\
        \bottomrule
    \end{tabular}
\end{table}

\subsection{Output Format}

Results are written to \texttt{intent\_classification.txt} as a TSV
file with columns \texttt{clause\_text}, \texttt{tfidf\_label},
\texttt{bert\_label}, and \texttt{confidence}.

\subsection{Comparative Example}

Table~\ref{tab:intent-example} compares the predictions of both models on
representative clauses.

\begin{table}[H]
    \centering
    \caption{Intent classification results for sample clauses.}
    \label{tab:intent-example}
    {\small
    \begin{tabularx}{\textwidth}{@{}>{\raggedright\arraybackslash}p{3.8cm}>{\raggedright\arraybackslash}p{2.0cm}>{\raggedright\arraybackslash}Xr@{}}
        \toprule
        \textbf{Clause Excerpt} & \textbf{TF-IDF} & \textbf{BERT} & \textbf{Conf.} \\
        \midrule
        \textit{``\ldots shall pay\ldots} & Obligation & Obligation & 0.96 \\
        \textit{``\ldots shall not disclose\ldots} & Prohibition & Prohibition & 0.98 \\
        \textit{``\ldots may terminate\ldots} & Right & Termination condition & 0.72 \\
        \textit{``\ldots terminates upon breach\ldots} & Termination condition & Termination condition & 0.95 \\
        \bottomrule
    \end{tabularx}
    }
\end{table}

The \textit{``may terminate''} clause illustrates a common ambiguity: TF-IDF's
n-gram features weight \textit{may} heavily toward the Right class, while BERT's
contextual representation captures the termination semantics of the full phrase
and predicts Termination Condition.  The lower confidence score (0.72) correctly
signals the model's uncertainty on this boundary case.

\ifSubfilesClassLoaded{
\begin{thebibliography}{1}

\bibitem{chalkidis2020legal}
I.~Chalkidis, M.~Fergadiotis, P.~Malakasiotis, N.~Aletras, and I.~Androutsopoulos,
\textit{LEGAL-BERT: The Muppets straight out of Law School},
in \textit{Findings of EMNLP 2020}, pp.\ 2898--2904, 2020.

\end{thebibliography}
}{}

\end{document}
